% Part: first-order-logic
% Chapter: proof-systems
% Section: seqeunt-calculus

\documentclass[../../../include/open-logic-section]{subfiles}

\begin{document}

\iftag{FOL}
      {\olfileid{fol}{prf}{seq}}
      {\olfileid{pl}{prf}{seq}}

\olsection{سیکوئنٹ حساب}

بھاویں بہت سارے !!{derivation} نظام !!{sentence}s دے بندوبستاں نال کم
کردے نیں، سیکوئنٹ حساب \emph{سیکوئنٹاں} نال کم کردا اے۔ سیکوئنٹ ایس
روپ دی اک عبارت اے
\[
!A_1, \dots, !A_m \Sequent !B_1, \dots, !B_m,
\]
یعنی !!{sentence}s دے سلسلیاں دی اک جوڑی، جہنوں سیکوئنٹ دی
علامت~$\Sequent$ وکھ کردی اے۔ دونیں سلسلیاں وچوں کوئی وی خالی ہو سکدا
اے۔ سیکوئنٹ حساب وچ !!^a{derivation} سیکوئنٹاں دا اک درخت ہُندا اے،
جتھے سب توں اُتّے والے سیکوئنٹ خاص روپ دے ہُندے نیں (اوہناں نوں
``ابتدائی سیکوئنٹ'' یا ``مسلّماں'' آکھیا جاندا اے) تے ہر ہور سیکوئنٹ
استنباط دے قاعدیاں وچوں کسے اک راہیں اپنے عین اُتّے والے سیکوئنٹاں توں
نکلدا اے۔ استنباط دے قاعدے یا تاں سیکوئنٹاں اندر !!{sentence}s نوں
بدلدے نیں (کھبے یا سجے پاسے اوہناں نوں شامل، خارج یا اگے پچھے کر کے)،
یا قاعدے دے نتیجے وچ اک مرکب !!{formula} متعارف کراؤندے نیں۔ مثال لئی،
$\LeftR{\land}$ دا قاعدہ $!A,
\Gamma \Sequent \Delta$ توں $!A \land !B, \Gamma \Sequent \Delta$ دا
استنباط کرن دی اجازت دیندا اے، تے $\RightR{\lif}$ دا قاعدہ $!A, \Gamma \Sequent
\Delta, !B$ توں $\Gamma \Sequent \Delta, !A \lif !B$ دا، ہر $\Gamma$،
$\Delta$، $!A$ تے~$!B$ لئی۔ (خاص طور تے، $\Gamma$ تے~$\Delta$ خالی
وی ہو سکدے نیں۔)

سیکوئنٹ حساب اُتے مبنی $\Proves$ تعلق دی تعریف ایویں کیتی جاندی اے:
$\Gamma \Proves !A$ اُدوں تے صرف اُدوں جدوں کوئی سلسلہ $\Gamma_0$ ایہو
جہا ہووے کہ ہر $!A$ جیہڑا $\Gamma_0$ وچ اے،~$\Gamma$ وچ ہووے، تے اک
!!{derivation} ہووے جہدی جڑ اُتے سیکوئنٹ~$\Gamma_0 \Sequent !A$ ہووے۔
$!A$ سیکوئنٹ حساب وچ اک قضیہ اے جے سیکوئنٹ~$\Sequent !A$ دا
!!a{derivation} ہووے۔ مثال لئی، ایہ !!a{derivation} وکھاؤندا اے کہ
$\Proves (!A \land !B) \lif !A$:
\begin{prooftree}
  \Axiom$!A \fCenter !A$
  \RightLabel{\LeftR{\land}}
  \UnaryInf$!A \land !B \fCenter !A$
  \RightLabel{\RightR{\lif}}
  \UnaryInf$\fCenter (!A \land !B) \lif !A$
\end{prooftree}

سیکوئنٹ حساب وچ سیٹ $\Gamma$ متضاد اے جے $\Gamma_0 \Sequent$ دا
!!a{derivation} ہووے (جتھے ہر $!A \in \Gamma_0$،~$\Gamma$
وچ ہووے تے سیکوئنٹ دا سجا پاسا خالی ہووے)۔ قاعدہ
\RightR{\Weakening} ورت کے، کسے متضاد سیٹ توں ہر !!{sentence} نوں
!!{derive} کیتا جا سکدا اے۔

سیکوئنٹ حساب نوں 1930 دی دہائی وچ Gerhard Gentzen نے ایجاد کیتا۔
اپنی منظم تے متناظر بناوٹ پاروں ایہ !!{derivation}s دا نظریہ تیار کرن
لئی بہت فائدہ مند رسمی ڈھانچا اے۔ سیکوئنٹ حساب وچ !!{derivation}s لبھنے
نسبتاً سَوکھے نیں، پر ایہ !!{derivation}s اکثر پڑھن وچ اوکھے ہُندے نیں
تے ثبوتاں نال اوہناں دا تعلق ویکھنا کدی کدی سَوکھا نہیں ہُندا۔ ایس دے باوجود
ایہ !!{derivation} نظاماں دا بہت سلیقے والا طریقہ ثابت ہویا اے، تے بہت
ساریاں منطقاں دے سیکوئنٹ حسابی نظام نیں۔

\end{document}
